%%%%%%%%%%%%%%%%%%%%%%%%%%%%%%%%%%%%%%%%%%%%%%%%%%%%%%%%%%%%%%%%%%%%%%%%%%%%%%%%%%
\begin{frame}[fragile]\frametitle{}
\begin{center}
{\Large Syntax}
\end{center}
\end{frame}

%%%%%%%%%%%%%%%%%%%%%%%%%%%%%%%%%%%%%%%%%%%%%%%%%%%%%%%%%%%%%%%%%%%%%%%%%%%%%%%%%%%
\begin{frame}[fragile]  \frametitle{What is Python?}
\begin{itemize}
\item Python is an interpreted, object-oriented, high-level programming language with dynamic semantics.
\item Python is simple and easy to learn, open source, free and cross-platform.
\item Python provides high-level built-in data structures.
\item Python is useful for rapid application development, as a scripting or glue language.
\item Python emphasizes readability, and supports modules and packages.
\end{itemize}
\end{frame}

%%%%%%%%%%%%%%%%%%%%%%%%%%%%%%%%%%%%%%%%%%%%%%%%%%%%%%%%%%%%%%%%%%%%%%%%%%%%%%%%%%%
\begin{frame}[fragile]  \frametitle{The Python Shell}
	\begin{itemize}
	\item Python can be run from a shell, IDE, or notebook -- commands are entered at the \texttt{>>>} prompt.
	\item Expressions are evaluated and the result is printed immediately.
	\item The prompt changes to \texttt{...} on continuation lines (loops, function definitions, etc.)
	\end{itemize}
\begin{lstlisting}
> python
Python 3.11.0 (default) [MSC v.1900 64 bit (AMD64)] on win32
>>> 2 + 2
4
>>> "hello" + \
...     " world!"
'hello world!'
\end{lstlisting}
\end{frame}

%%%%%%%%%%%%%%%%%%%%%%%%%%%%%%%%%%%%%%%%%%%%%%%%%%%%%%%%%%%%%%%%%%%%%%%%%%%%%%%%%%%
\begin{frame}[fragile] \frametitle{Lines}
\begin{itemize}
\item Statement separator: a semicolon.
\item Only needed for multiple statements on one line -- otherwise, nothing.
\item Continuation: a backslash.
\item But not necessary inside a ``context'': \texttt{( )}, \texttt{\{ \}}, \texttt{[ ]}.
\end{itemize}
\end{frame}

%%%%%%%%%%%%%%%%%%%%%%%%%%%%%%%%%%%%%%%%%%%%%%%%%%%%%%%%%%%%%%%%%%%%%%%%%%%%%%%%%%%
\begin{frame}[fragile] \frametitle{Blocks and Indentation}
\begin{itemize}
\item Blocks are formed by indentation -- no begin/end brackets.
\item Indentation is 4 spaces, no hard tabs.
\item \lstinline{if, while, def, except}, etc. end with a colon.
\item A colon must be followed by an indented block on the next line.
\item One script can be run directly, or imported by other modules.
\end{itemize}
\end{frame}

%%%%%%%%%%%%%%%%%%%%%%%%%%%%%%%%%%%%%%%%%%%%%%%%%%%%%%%%%%%%%%%%%%%%%%%%%%%%%%%%%%%
\begin{frame}[fragile] \frametitle{Comments}
Everything after \texttt{\#} is ignored -- no block comments in Python.
\begin{lstlisting}
# compute the percentage of the hour that has elapsed
percentage = (minute * 100) / 60

v = 5             # useless/redundant comment
v = 5             # velocity in meters/second -- useful comment
\end{lstlisting}
Use a comment to temporarily ``disable'' a line of code, rather than deleting it.
\end{frame}

%%%%%%%%%%%%%%%%%%%%%%%%%%%%%%%%%%%%%%%%%%%%%%%%%%%%%%%%%%%%%%%%%%%%%%%%%%%%%%%%%%%
\begin{frame}[fragile]\frametitle{f-Strings: Modern String Formatting (Python 3.6+)}
Embed variables and expressions directly in strings by prefixing with \texttt{f}. f-strings are \textbf{faster}, \textbf{more readable}, and the \textbf{recommended} style for all new Python 3 code.
\begin{lstlisting}
name = "Alice"
age  = 30

# Old %-style (avoid in new code)
print("Hello, %s. You are %d." % (name, age))
# .format() style
print("Hello, {}. You are {}.".format(name, age))
# f-string (preferred since Python 3.6)
print(f"Hello, {name}. You are {age}.")
print(f"Next year: {age + 1}")        # expressions work too
print(f"Pi: {3.14159:.4f}")           # format specs work too
\end{lstlisting}
\end{frame}

%%%%%%%%%%%%%%%%%%%%%%%%%%%%%%%%%%%%%%%%%%%%%%%%%%%%%%%%%%%%%%%%%%%%%%%%%%%%%%%%%%%
\begin{frame}[fragile]  \frametitle{Assignment}
	\begin{itemize}
	\item Assignment creates references, not values: \lstinline{tmp = "hello"; tmp = 10} -- the string gets deallocated.
	\item Augmented assignment is valid: \lstinline{x += 1}.
	\item Pre/post increment/decrement (\lstinline{x++}, \lstinline{++x}) are \textbf{not} valid Python -- there is no \texttt{++} operator.
	\item Multiple assignment to one object: \lstinline{x = y = z = 1}.
	\item Tuple unpacking: \lstinline{(x, y, z) = (3.5, 5.5, 'string')}.
	\item Swapping values: \lstinline{(x, y) = (y, x)}.
	\end{itemize}
\end{frame}

%%%%%%%%%%%%%%%%%%%%%%%%%%%%%%%%%%%%%%%%%%%%%%%%%%%%%%%%%%%%%%%%%%%%%%%%%%%%%%%%%%%
\begin{frame}[fragile]\frametitle{Built-in Object Types}
	\begin{itemize}
	\item Numbers: \lstinline{3.1415, 1234, 999, 3+4j}
	\item Strings: \lstinline{'spam', "guido's"}
	\item Lists: \lstinline{[1, [2, 'three'], 4]}
	\item Dictionaries: \lstinline|{'food': 'spam', 'taste': 'yum'}|
	\item Tuples: \lstinline{(1, 'spam', 4, 'U')}
	\item Sets: \lstinline|{1, 2, 3, 'foo', 'bar'}|
	\end{itemize}
\end{frame}

%%%%%%%%%%%%%%%%%%%%%%%%%%%%%%%%%%%%%%%%%%%%%%%%%%%%%%%%%%%%%%%%%%%%%%%%%%%%%%%%%%%
\begin{frame}[fragile]\frametitle{Numbers}
	\begin{itemize}
	\item Integers (unlimited precision): \lstinline{1234, -24, 0, 999999999999}
	\item Float: \lstinline{3.1415, 2.7122}
	\item Octal and hex: \lstinline{0o177, 0x9ff}
	\item Complex: \lstinline{3+4j, 3.0+4.0j, 3J}
	\end{itemize}
\end{frame}

%%%%%%%%%%%%%%%%%%%%%%%%%%%%%%%%%%%%%%%%%%%%%%%%%%%%%%%%%%%%%%%%%%%%%%%%%%%%%%%%%%%
\begin{frame}[fragile]\frametitle{Strings (immutable sequences)}
	\begin{itemize}
	\item single quote \lstinline{s1 = 'egg'}, double quotes \lstinline{s2 = "spam's"}, triple quotes \lstinline{block = """..."""}
	\item concatenate \lstinline{s1 + s2}, repeat \lstinline{s2 * 3}
	\item index, slice \lstinline{s2[i], s2[i:j]}, length \lstinline{len(s2)}
	\item formatting \lstinline|"a {} parrot".format('dead')|
	\item iteration \lstinline{for x in s2}, membership \lstinline{'m' in s2}
	\end{itemize}
\end{frame}

%%%%%%%%%%%%%%%%%%%%%%%%%%%%%%%%%%%%%%%%%%%%%%%%%%%%%%%%%%%%%%%%%%%%%%%%%%%%%%%%%%%
\begin{frame}[fragile]\frametitle{Lists}
	\begin{itemize}
	\item Ordered collections of arbitrary objects, accessed by offset.
	\item Variable length, heterogeneous, arbitrarily nest-able.
	\item Mutable -- an array of object references.
	\end{itemize}
\begin{lstlisting}
L = []                       # empty list
L2 = [0, 1, 2, 3]             # four items
L3 = ['abc', ['def', 'ghi']]  # nested
L2[i], L3[i][j]               # index
L2[i:j], len(L2)              # slice, length
L1 + L2, L2 * 3                # concatenate, repeat
L2.append(4); L2.sort(); L2.reverse()
del L2[k]; L2[i] = 1; L2[i:j] = [4, 5, 6]
\end{lstlisting}
\end{frame}

%%%%%%%%%%%%%%%%%%%%%%%%%%%%%%%%%%%%%%%%%%%%%%%%%%%%%%%%%%%%%%%%%%%%%%%%%%%%%%%%%%%
\begin{frame}[fragile]\frametitle{Dictionaries}
	\begin{itemize}
	\item Accessed by key, not offset -- unordered collections of arbitrary objects.
	\item Variable length, heterogeneous, arbitrarily nest-able, mutable.
	\item Implemented as hash tables of object references.
	\end{itemize}
\begin{lstlisting}
d1 = {}                              # empty
d2 = {'spam': 2, 'eggs': 3}          # two items
d3 = {'food': {'ham': 1, 'egg': 2}}  # nested
d2['eggs'], d3['food']['ham']        # indexing
d2.keys(), d2.values(), len(d1)
d2[key] = new                        # add/change
del d2[key]                          # delete
\end{lstlisting}
\end{frame}

%%%%%%%%%%%%%%%%%%%%%%%%%%%%%%%%%%%%%%%%%%%%%%%%%%%%%%%%%%%%%%%%%%%%%%%%%%%%%%%%%%%
\begin{frame}[fragile]\frametitle{Tuples and Files}
	\begin{itemize}
	\item \textbf{Tuples} are like lists but immutable -- use them when the content shouldn't change.
	\item \textbf{Files}:
	\end{itemize}
\begin{lstlisting}
f = open('data', 'r')
s = f.read()          # read all
s = f.read(n)          # read n bytes
s = f.readline()       # read next line
lines = f.readlines()  # read into a list
out = open('/tmp/spam', 'w')
out.write(s); out.writelines(lines); out.close()
\end{lstlisting}
\end{frame}

% %%%%%%%%%%%%%%%%%%%%%%%%%%%%%%%%%%%%%%%%%%%%%%%%%%%%%%%%%%%%%%%%%%%%%%%%%%%%%%%%%%%
% \begin{frame}[fragile]\frametitle{Comparisons vs. Equality}
% \begin{lstlisting}
% L1 = [1, ('a', 3)]
% L2 = [1, ('a', 3)]
% L1 == L2   # True -- == tests value equivalence
% L1 is L2   # False -- is tests object identity
% \end{lstlisting}
% \end{frame}

%%%%%%%%%%%%%%%%%%%%%%%%%%%%%%%%%%%%%%%%%%%%%%%%%%%%%%%%%%%%%%%%%%%%%%%%%%%%%%%%%%%
\begin{frame}[fragile]\frametitle{Control Flow: if, while, for}
\begin{lstlisting}
if not done and (x > 1):
    doit()
elif done and (x <= 1):
    dothis()
else:
    dothat()

while True:
    line = read_line()
    if len(line) == 0:
        break

for letter in 'hello world':
    print(letter)
for i in range(2, 10, 2):
    print(i)
\end{lstlisting}
\end{frame}

% %%%%%%%%%%%%%%%%%%%%%%%%%%%%%%%%%%%%%%%%%%%%%%%%%%%%%%%%%%%%%%%%%%%%%%%%%%%%%%%%%%%
% \begin{frame}[fragile]\frametitle{pass}
% Temporary filler, the stub. Wherever a colon requires an indented block but there's nothing to do yet, put \lstinline{pass} on the next line.
% \begin{lstlisting}
% def not_yet_implemented():
%     pass
% \end{lstlisting}
% \end{frame}

% %%%%%%%%%%%%%%%%%%%%%%%%%%%%%%%%%%%%%%%%%%%%%%%%%%%%%%%%%%%%%%%%%%%%%%%%%%%%%%%%%%%
% \begin{frame}[fragile]\frametitle{Errors and Exceptions}
% 	\begin{itemize}
% 	\item \texttt{NameError}: undeclared variable. \texttt{ZeroDivisionError}: division by zero.
% 	\item \texttt{IndexError}: out-of-range sequence index. \texttt{KeyError}: missing dictionary key.
% 	\item \texttt{OSError}: input/output error. \texttt{AttributeError}: unknown object attribute.
% 	\end{itemize}
% \begin{lstlisting}
% try:
%     f = open('blah')
% except OSError:
%     print('could not open file')
% \end{lstlisting}
% \end{frame}

%%%%%%%%%%%%%%%%%%%%%%%%%%%%%%%%%%%%%%%%%%%%%%%%%%%%%%%%%%%%%%%%%%%%%%%%%%%%%%%%%%%
\begin{frame}[fragile]\frametitle{Functions}
	\begin{itemize}
	\item Functions can return any type of object -- \lstinline{None} is returned by default.
	\item Multiple values can be returned; parameters can have default arguments.
	\item Variable-length arguments and anonymous (\lstinline{lambda}) functions are supported.
	\end{itemize}
\begin{lstlisting}
def test(a, b=2, d=some_func):
    return d(a, b)

test(3)
test(b=4, a=3)
test(1, 2, lambda x, y: x * y)
\end{lstlisting}
\end{frame}

% %%%%%%%%%%%%%%%%%%%%%%%%%%%%%%%%%%%%%%%%%%%%%%%%%%%%%%%%%%%%%%%%%%%%%%%%%%%%%%%%%%%
% \begin{frame}[fragile]\frametitle{Modules, Namespaces and Packages}
% 	\begin{itemize}
% 	\item A file is a module, e.g. \texttt{myio.py}, containing a function \texttt{load}.
% 	\item Code in \texttt{myio.py} lives in the \texttt{myio} namespace.
% 	\item Packages are bundles of modules (a directory with an \texttt{\_\_init\_\_.py}).
% 	\end{itemize}
% \begin{lstlisting}
% import myio
% myio.load()
%
% from myio import load
% load()
% \end{lstlisting}
% \end{frame}

% %%%%%%%%%%%%%%%%%%%%%%%%%%%%%%%%%%%%%%%%%%%%%%%%%%%%%%%%%%%%%%%%%%%%%%%%%%%%%%%%%%%
% \begin{frame}[fragile]\frametitle{Class}
% \begin{lstlisting}
% class Cone:
%     def __init__(self, d0, de, L):
%         self.a0 = d0 / 2
%         self.ae = de / 2
%         self.L = L
%
%     def radius(self, z):
%         return self.ae + (self.a0 - self.ae) * z / self.L
%
% c = Cone(0.1, 0.2, 1.5)
% c.radius(0.5)
% \end{lstlisting}
% \end{frame}

% %%%%%%%%%%%%%%%%%%%%%%%%%%%%%%%%%%%%%%%%%%%%%%%%%%%%%%%%%%%%%%%%%%%%%%%%%%%%%%%%%%%
% \begin{frame}[fragile]\frametitle{Standard Library Modules}
% 	\begin{itemize}
% 	\item \textbf{os}: file and process operations.
% 	\item \textbf{time}: dates and times.
% 	\item \textbf{string} / \textbf{re}: string operations and regular expressions.
% 	\item \textbf{copy}: object copying.
% 	\item \textbf{NumPy}: numerical array processing (third-party, not stdlib, but ubiquitous).
% 	\item Visit \url{https://pypi.org} for the full package ecosystem.
% 	\end{itemize}
% \end{frame}

%%%%%%%%%%%%%%%%%%%%%%%%%%%%%%%%%%%%%%%%%%%%%%%%%%%%%%%%%%%%%%%%%%%%%%%%%%%%%%%%%%%
\begin{frame}[fragile]\frametitle{Quick Check: Syntax}
\begin{block}{Think About It}
C and Java use braces \{ \} to mark blocks -- indentation is just a human-readable convention on top of that. Python uses indentation itself as the block boundary. What's the tradeoff?
\end{block}
\pause
\begin{block}{Answer}
It forces what the code \emph{looks} like and what it \emph{does} to always agree -- you can't have a block that visually appears nested one way but executes another, a class of bug that misleading brace placement can cause in C or Java. The tradeoff is that whitespace becomes meaningful: mixing tabs and spaces, or one stray space, can silently change what the code does, without an obvious visual cue like a misplaced brace.
\end{block}
\end{frame}
